% SPDX-License-Identifier: CC-BY-SA-4.0
% Author: Matthieu Perrin
% Part: <Nom de la partie>
% Section: <Nom de la section>
% Sub-section: <Nom de la sous-section>  % (facultatif, laisser vide si non utilisé)
% Frame: <Titre de la slide>

\begingroup

\begin{frame}{Le problème universel borné}

  \on[text, top=-4mm]{
    \Probleme{NPUniversal}{
      \begin{itemize}
      \item un \structure{encodage} $\mathit{encode}(M)$ d'une \alert{MTND} $M$,
      \item un mot $u$
      \item un entier $k$ \structure{écrit en unaire}
      \end{itemize}
    }{
      Est-ce que $M$ accepte $u$ en moins de $k$ étapes ?
    }
  }

  \onBlock[bottom=-4mm]{Théorème -- Complexité de \textsc{NPUniversal}}{
    Le problème \textsc{NPUniversal} est \textsc{np}-complet.

    \vspace{2mm}
    \structure{$\textsc{NPUniversal} \in \textsc{np}$.} On adapte la machine universelle (page~\ref{limits/universality/semidecidable/mtu}) 
    \begin{itemize}
    \item Choix non-déterministe d'une transition de $M$ 
    \item Décompte des transitions de $M$ borné par $k$ 
    \end{itemize}

    \vspace{1mm}
    \structure{$\textsc{NPUniversal} \in \textsc{np-difficile}$.} Soit $A \in \textsc{np}$. Montrons que $A \leq_p \textsc{NPUniversal}$.
    \begin{itemize}
    \item Il existe $M_A$ une MTND qui accepte $A$ en au plus $P(n)$ étapes.
    \item On a la réduction polynomiale $\varphi$ de $A$ vers $\textsc{NPUniversal}$ suivante :

      \vspace{-2mm}
      $$\alert{\varphi : \left\{\begin{array}{rcl}
        \mathcal{I}(A) & \rightarrow & \mathcal{I}(\textsc{NPUniversal}) \\
        u              & \mapsto     & \langle \mathit{encode}(M_A), u, 1^{P(|u|)} \rangle \\
        \end{array}\right.}$$ 
    \end{itemize}
  }

\end{frame}

\endgroup
